\section{BHSM-native transport of the physical topographic mode}
\label{sec:bhsm-native-transport}

The cosmological bridge is formulated first in the mode ontology of the
Berger--Hopf Standard Model (BHSM), and only secondarily in a four-dimensional
effective-field-theory coordinate system.  A topographic carrier is therefore
not introduced as an additional free bulk scalar.  It is an action-owned
geometric/boundary mode whose physical dynamics are assigned after the
constraint and gauge quotients have been taken.

\subsection{Physical tangent space and reduced operator}

At a background configuration $z_\star$, define
\begin{equation}
T_{\rm phys}(z_\star)
=
\ker D\mathcal C(z_\star)\cap\mathcal G^\perp ,
\end{equation}
where $\mathcal C$ is the action-owned constraint map and $\mathcal G$ is the
space of infinitesimal gauge directions.  Let
\begin{equation}
\widehat{\mathcal H}
=
P_{\rm phys}
D^2\mathcal L_{\rm KKT}(z_\star)
P_{\rm phys}
\end{equation}
be the physical Hessian.  For the first admissible scalar-topographic harmonic
projector $P_2$, define the reduced operator
\begin{equation}
\mathcal K_2
=
P_2
\left[
\widehat{\mathcal H}
-
\widehat{\mathcal H}Q_2
\left(
Q_2\widehat{\mathcal H}Q_2
\right)^{-1}
Q_2\widehat{\mathcal H}
\right]
P_2 ,
\label{eq:bhsm-K2}
\end{equation}
where $Q_2=1-P_2$ on the already physical tangent space and the inverse is
taken only on the admissible eliminated block.

\subsection{Action-owned $n=2$ envelope transport theorem}

\paragraph{Theorem 1.}
Assume that the cosmological parent envelope is
$\mathcal E_{\rm cos}\simeq S^3(R_H)$, that scalar/topographic perturbations
are reduced to $T_{\rm phys}$ before physical propagation is assigned, that
$n=0$ is absorbed into the homogeneous background, and that the exceptional
regular closed-FLRW scalar $n=1$ sector is removed by the physical quotient.
If the retained $n=2$ kinetic form is positive and nondegenerate, then the
first physical scalar-topographic carrier is an action-owned $n=2$ envelope
mode.  Its canonical momentum is generated by the same reduced action and is
not a separately postulated bulk field.

\paragraph{Proof.}
Let $e_2(x)$ be a normalized element of $P_2T_{\rm phys}$ and write
\begin{equation}
\delta\Phi_{\rm topo}(t,x)=q_2(t)e_2(x).
\end{equation}
The gauge- and constraint-reduced quadratic action has the one-mode form
\begin{equation}
S^{(2)}_2
=
\frac12\int dt\,a^3
\left[
\mathcal G_2(t)(D_tq_2)^2
-
\mathcal M_2^2(t)q_2^2
+
2\mathcal J_2(t)q_2
\right],
\label{eq:bhsm-q2-action}
\end{equation}
with coefficients obtained from matrix elements of the reduced physical
operator.  The conjugate momentum is
\begin{equation}
\Pi_2=a^3\mathcal G_2D_tq_2 ,
\end{equation}
and variation gives
\begin{equation}
D_t\!\left(a^3\mathcal G_2D_tq_2\right)
+
a^3\mathcal M_2^2q_2
=
a^3\mathcal J_2 .
\label{eq:bhsm-q2-eom}
\end{equation}
Thus $(q_2,\Pi_2)$ is the reduced phase-space pair of one action-owned
geometric mode. \hfill$\square$

\subsection{Effective closed-Horndeski matching}

The exact closed-Horndeski $n=2$ result derived elsewhere in this manuscript
is used as an effective kinetic representation of the same retained mode,
rather than as its microscopic ontology.  In the Bellini--Sawicki convention,
\begin{equation}
\mathcal G_{S,2}
=
M_{\rm Pl}^2
\frac{
5D_{\rm kin}
}{
D_{\rm kin}
+
10\left(1-\alpha_B^{\rm BS}/2\right)^2
}.
\end{equation}
Hence $D_{\rm kin}>0$ and a nonsingular denominator provide the effective
positive-kinetic matching condition used on the reference branch.  The EFT
Stueckelberg coordinate $\pi$ is consequently interpreted as a coordinate
representative of the retained response of $q_2$, not as an additional BHSM
degree of freedom.

\subsection{Exact observer dipole}

A degree-two scalar harmonic on $S^3$ can be written
\begin{equation}
T_2=A_{AB}X^AX^B,\qquad A_A{}^A=0 .
\end{equation}
On an observer shell,
\begin{equation}
T_2
=
A_{44}\left(c^2-\frac{s^2}{3}\right)
+
2cs\,\mathbf a\!\cdot\!\hat{\mathbf n}
+
s^2S_{ij}\hat n^i\hat n^j .
\end{equation}
The pure $A_{4i}$ subspace therefore gives the exact sky dipole
\begin{equation}
\boxed{
T_{2,\rm dip}
=
A\sin(2\psi)
\left(
\hat{\mathbf n}\!\cdot\!\hat{\mathbf p}
\right)
}.
\label{eq:bhsm-n2-exact-dipole}
\end{equation}

\subsection{Curved-EFT compatibility status}

The curvature-aware EFT system remains an important representation-level
cross-check.  It is not used to define the existence or mode count of the
BHSM carrier.  The literal curved constraint/evolution equations and the
reference background have been transcription-audited, while the global
curved-EFT differential-algebraic implementation retains a nonzero
constraint-preservation defect.  That unresolved object is therefore
classified as an EFT compatibility audit rather than as the foundational
transport law.  The defining bridge evolution is Eq.~\eqref{eq:bhsm-q2-eom}.
